\section{A diagnostic for projection onto the full efficient core}
\label{net:projection}
The following examples distinguish the full efficient core from a smaller
set of dual-price-generated allocations.  They concern the specific rule
that projects grand-minus-complement marginal payments onto that full core.
They are source-comparison diagnostics, not new general subsidy lower bounds.

For four fixed owners $A=\{0,1\}$, $B=\{2,3\}$, initially put
\[
 v_c(S)=\min\left\{\sum_{i\in S\cap A}c_i,
                         \sum_{j\in S\cap B}c_j\right\},\qquad
 \mathcal C(c)=\{x:x(I)=T,\ x(S)\ge v_c(S)\text{ for every }S\}.
\]
Singleton constraints imply nonnegativity.  Define
$U_i(c)=T-v_c(I\setminus\{i\})$ and let $P(c)$ be the unique minimizer
of $\frac12\|x-U(c)\|_2^2$ over $\mathcal C(c)$.
Compare the profiles
\[
 L=(1/2,1,1,1),\qquad H=(1/2,1,2,1),
\]
which differ only in owner 2's own capacity.

\begin{proposition}[Linear full-core projection failure]
\label{net:linear-projection}
The stated projection pays $P_2(L)=1/4$ and $P_2(H)=0$.
These games arise from four unique resource types and four fixed linear
production activities.  Their dual-price allocation set, however, is already
a singleton at the low profile, so this example does not refute projection
onto that smaller set.
\end{proposition}
\begin{proof}
At $L$, $T=3/2$ and $U=(1/2,1,1/2,1/2)$.
Orthogonal projection onto the efficiency hyperplane subtracts $1/4$ from
every coordinate, giving
\[
 x^L=(1/4,3/4,1/4,1/4).
\]
It lies in the full core.  Empty and single-side coalitions have value zero.
The mixed-pair values are $1/2,1/2,1,1$, equal to their payments.
The triples $\{0,1,j\}$ have value one and payment $5/4$;
$\{0,2,3\}$ has value $1/2$ and payment $3/4$;
$\{1,2,3\}$ has value one and payment $5/4$.
The grand value and payment are $3/2$.
Thus the hyperplane minimizer is feasible for the core and is its unique
projection by strict convexity.

At $H$, the marginal vector is $(1/2,1,1/2,0)$, not the core allocation.
The coalition $\{0,1,2\}$ has value $3/2$, so efficiency and singleton
coverage force $x_3=0$ in every core allocation.  The coalitions
$\{0,3\}$ and $\{1,3\}$ then give $x_0\ge1/2$ and $x_1\ge1$;
efficiency and nonnegativity force $x=(1/2,1,0,0)$.
This vector covers every coalition because the sum of its present $A$-side
capacities bounds production from above.  It is therefore the singleton high
core, proving the payment drop.

For an explicit production representation, activity $y_{ij}$ consumes one
unit of each resource $i\in A,j\in B$ and earns one unit of profit.
The four constraints are
\[
 y_{02}+y_{03}\le c_0,\quad y_{12}+y_{13}\le c_1,\quad
 y_{02}+y_{12}\le c_2,\quad y_{03}+y_{13}\le c_3.
\]
The maximum total activity equals the minimum of the two pool totals:
each total is an upper bound, attained by a bipartite transportation flow.
Giving nonmembers capacity zero realizes exactly the characteristic game.

At both profiles the unique optimal nonnegative resource-price vector is
$(1,1,0,0)$.  Indeed, let $a=\min(p_0,p_1)$ and $b=\min(p_2,p_3)$.
Dual feasibility $p_i+p_j\ge1$ implies $a+b\ge1$.
Writing the pool totals as $A_0=3/2$ and $B_0>A_0$, the objective is at least
$A_0a+B_0b\ge A_0+(B_0-A_0)b$.
Equality at the optimum forces $b=0,a=1$ and, because all capacities are
positive, forces both prices within each pool to equal their minima.
Thus dual-price-generated owner payments are $(1/2,1,0,0)$ even at $L$.
This is strictly smaller than the low full core used by the diagnostic.
\end{proof}

\begin{proposition}[Strictly concave utility does not remove the failure]
\label{net:concave-projection}
Keep the same four owners, activities, resource constraints, and reports,
but maximize $\sum_{i\in A,j\in B}u(y_{ij})$ with
$u(t)=10t/(10+t)$.  Projection of the corresponding marginal vector onto
the full efficient core still violates own-report monotonicity: owner 2's
payment falls by at least $1270/9471$ between $L$ and $H$.
This nonlinear production value is not asserted to be homogeneous.
\end{proposition}
\begin{proof}
Here $u(0)=0$, $u'(t)=100/(10+t)^2>0$, and
$u''(t)=-200/(10+t)^3<0$ on the nonnegative half-line.
Set
\[
 a=2u(1/4)=20/41,\qquad b=2u(1/2)=20/21,\qquad
 T=a+b=1240/861.
\]
At both profiles split each $A$ row equally between the two $B$ columns;
each column then uses $3/4$.  The split is feasible and attains the upper
bound from row-wise concavity, so it is optimal.
At $L$, deleting owners $0,1,2,3$ gives values $(b,a,b,b)$.
For example, with only one $B$ column of capacity one, two rows split its
capacity equally, which is feasible and optimal by concavity.
Thus $U(L)=(a,b,a,a)$ and its efficiency-hyperplane projection is
\[
 x^L=(a/2,b-a/2,a/2,a/2).
\]
It is nonnegative.  Mixed pairs involving owner 0 have value $u(1/2)\le a$
and payment $a$; those involving owner 1 have value $u(1)\le b$ and payment
$b$.  Triples $\{0,1,j\}$ have value $b$ and payment $b+a/2$;
$\{0,2,3\}$ has value $a$ and payment $3a/2$;
$\{1,2,3\}$ has value $b$ and payment $b+a/2$.
Other proper coalitions have value zero.  Efficiency holds, so $x^L$ is the
unique full-core projection.

At $H$, coalition $\{0,1,2\}$ has value $u(1/2)+u(1)$, because owner 2
can accommodate both full rows.  Put
\[
 \delta=T-u(1/2)-u(1)=\frac{520}{9471}>0.
\]
For every high core allocation, $x_3\le\delta$.
The coalitions $\{0,3\}$ and $\{1,3\}$ give
$x_0+x_3\ge u(1/2)$ and $x_1+x_3\ge u(1)$.
Summing and using efficiency implies
$x_2\le\delta+x_3\le2\delta=1040/9471$.
The high core is nonempty: $(a,b,0,0)$ covers every coalition, since its
present $A$ rows contribute at most their row-concavity bounds $a,b$.
The unique high projection therefore exists and satisfies the same bound,
regardless of its other coordinates.  Owner 2's low payment is
$a/2=10/41$, so its decrease is at least
$10/41-2\delta=1270/9471>0$.
\end{proof}

\paragraph{Source-applicability limit.}
The inspected passages of Amigo, Belzarena, and Vaton (2016),
\emph{Revenue sharing in network utility maximization problems}
(\href{https://doi.org/10.1007/s11066-016-9112-4}{doi:10.1007/s11066-016-9112-4}),
describe ordinary full-core constraints and Euclidean
projection of grand-minus-complement marginals.  The assumptions preceding its
Problem~1 have not all been recovered.  Accordingly, applicability of these
examples to that source's complete model remains conditional.  The strictly
concave example removes one possible utility mismatch, but it does not establish
all of the source's hypotheses.  The inspected conclusion reports simulation
evidence for monotonicity; these examples are not presented as a refutation
of a proved all-profile theorem.  Both propositions concern exactly the
projection target defined here and do not apply to an unspecified restricted
allocation set or to the main homogeneous-production theorem.
